\chapter{Nasal Swab Culture}

\section*{What Is This Test?}
A nasal swab culture involves collecting a specimen from the anterior nares or nasopharynx for microbiologic identification of colonizing or infecting organisms. The test is most commonly performed to screen for nasal carriage of \textit{Staphylococcus aureus} (including methicillin-resistant \textit{S. aureus} [MRSA]) as part of infection control programs or preoperative screening. Nasal swabs are also used for detection of respiratory viruses (influenza, SARS-CoV-2, RSV) by antigen or molecular testing, for pertussis diagnosis (nasopharyngeal swab for \textit{Bordetella pertussis} PCR or culture), and for screening of multidrug-resistant organisms (MDROs) such as vancomycin-resistant enterococci (VRE) and carbapenem-resistant Enterobacterales (CRE) in selected patient populations. The anterior nares are the primary reservoir for \textit{S. aureus}: approximately 20--30\% of healthy adults are persistent nasal carriers, 30\% are intermittent carriers, and approximately 50\% are non-carriers. Nasal carriage of \textit{S. aureus} is the major risk factor for subsequent \textit{S. aureus} infection in surgical patients, patients on hemodialysis, and ICU patients \cite{wertheim2005role}. Active detection and isolation (ADI) programs screening high-risk patients on admission for MRSA nasal carriage are standard in many healthcare systems.

\section*{Alternative Names}
Nasal Culture; Nasopharyngeal Swab; NP Swab; MRSA Nasal Screen; Anterior Nares Culture; Nasal Swab for Staph; Nasal Surveillance Culture

\section*{Principle}
Nasal swab consists of specimen collection using a flocked nylon or rayon-tipped swab from the anterior nares or nasopharynx, followed by microbiologic analysis. For MRSA screening, the swab is inoculated onto selective chromogenic agar (e.g., CHROMagar MRSA, BBL CHROMagar Staph aureus) containing cefoxitin or oxacillin, incubated at 35--37\textdegree{}C for 18--24 hours, and examined for characteristic colony color (mauve/pink). Alternatively, molecular methods using real-time PCR detect the \textit{mecA} or \textit{mecC} gene (conferring methicillin resistance) and \textit{S. aureus}-specific targets (e.g., \textit{nuc}, \textit{spa}) directly from the swab, providing results in 1--2 hours (Cepheid GeneXpert MRSA, Roche Cobas MRSA/SA, BD MAX MRSA). For respiratory viral detection, the swab is placed in viral transport medium and tested by RT-PCR, direct fluorescent antibody (DFA), or rapid antigen detection. For \textit{Bordetella pertussis}, the swab is plated on Regan-Lowe selective agar or tested by PCR targeting insertion sequence IS481. For bacterial culture, swabs are plated on blood agar, MacConkey agar, and selective media as indicated, with identification by MALDI-TOF mass spectrometry or biochemical methods and antimicrobial susceptibility testing by broth microdilution or disk diffusion.

\section*{History}
The practice of nasal swab culture for \textit{S. aureus} carriage dates to the 1930s, when the epidemiology of staphylococcal infections in hospitals began to be elucidated. The emergence of MRSA in the 1960s spurred development of selective media and screening programs. The first chromogenic agar for MRSA detection was introduced in the 1990s, significantly reducing the time to detection compared to conventional culture. Molecular detection of MRSA from nasal swabs by PCR became available in the mid-2000s, with the Cepheid GeneXpert MRSA test receiving FDA clearance in 2007. The superiority of molecular methods for rapid MRSA screening was demonstrated in a landmark study (Robotham et al., 2011) showing that PCR-based screening reduced MRSA transmission in surgical and ICU patients by allowing pre-emptive contact isolation. The COVID-19 pandemic in 2020 dramatically expanded the use of nasal and nasopharyngeal swabs for widespread RT-PCR testing, standardizing the collection technique globally. Current guidelines from SHEA, IDSA, and APIC recommend active surveillance cultures for MRSA in high-risk populations, though the universal versus targeted approach remains debated \cite{calfee2014strategies}.

\section*{Why This Test Is Ordered}
\begin{itemize}
  \item Screening for nasal carriage of MRSA in patients admitted to the ICU, undergoing major surgery (cardiothoracic, orthopedic with implant, neurosurgery), on hemodialysis, or transferred from long-term care facilities, to guide contact precautions and decolonization \cite{safdar2006effectiveness}
  \item Preoperative \textit{S. aureus} screening prior to high-risk surgical procedures (total joint arthroplasty, cardiac surgery) to identify carriers who benefit from preoperative decolonization with intranasal mupirocin and chlorhexidine bathing, which reduces surgical site infections by 40--60\% \cite{bode2010preventing}
  \item Diagnosis of respiratory viral infections (influenza, SARS-CoV-2, RSV) by antigen or molecular testing of nasopharyngeal swab specimens during respiratory illness season
  \item Diagnosis of pertussis (\textit{Bordetella pertussis}) by nasopharyngeal PCR or culture in patients with prolonged paroxysmal cough, post-tussive emesis, or whoop
  \item Investigation of nosocomial outbreaks of MRSA or other MDROs to identify colonized healthcare workers or environmental sources
  \item Screening for nasal \textit{S. aureus} colonization in patients with recurrent staphylococcal skin and soft tissue infections to determine if a decolonization strategy is indicated
\end{itemize}

\section*{Primary vs Secondary Test}
Nasal swab for MRSA screening is a primary surveillance test indicated by institutional infection control policy, not by clinical symptoms. For diagnostic purposes (respiratory virus detection, pertussis, bacterial sinusitis), the nasal swab serves as a primary specimen collection method for laboratory testing.

\section*{How This Test Is Performed}
For anterior nares MRSA screening: A flocked swab is moistened with sterile saline, inserted 1--2 cm into one anterior nare, rotated against the nasal mucosa for 5 rotations, then the same swab is used to sample the other nare. For nasopharyngeal (NP) collection for respiratory pathogens: A flexible flocked swab is inserted into one nostril, advanced along the floor of the nasal cavity to the posterior nasopharynx (distance approximately equal to that from the nose to the external ear), rotated against the mucosa for 5--10 seconds, then withdrawn and placed into transport medium. For combined MRSA and respiratory pathogen detection, two separate swabs (anterior nares and nasopharyngeal) are collected. Specimens for bacterial culture should be transported at room temperature in Amies transport medium and processed within 24 hours. Specimens for molecular testing are typically collected with manufacturer-specific dry swabs or swabs placed in the assay-specific transport buffer.

\section*{What Preparation Is Needed}
No special patient preparation is required. Avoid intranasal antibiotic ointments (e.g., mupirocin) for at least 48 hours before specimen collection if assessing MRSA carriage or following decolonization protocols. Avoid eating, drinking, or tooth brushing for 30 minutes before nasopharyngeal specimen collection for molecular respiratory testing to prevent specimen contamination.

\section*{Contraindications and Precautions}
\begin{itemize}
  \item Nasopharyngeal swab insertion is relatively contraindicated in patients with severe coagulopathy, thrombocytopenia, recent nasal surgery, or severe nasal obstruction (deviated septum, nasal polyps). Gag reflex or coughing during NP specimen collection may generate aerosols; N95/FFP2 respirators and eye protection are indicated when collecting from patients with suspected highly transmissible respiratory pathogens. False-negative MRSA culture may occur if the patient has recently received antibiotics active against MRSA (vancomycin, linezolid) or if specimens are collected soon after chlorhexidine bathing. MRSA PCR detects the \textit{mecA} gene but does not distinguish viable from nonviable organisms. \textit{S. aureus} isolates lacking \textit{mecA} but with borderline oxacillin MIC (2--4 mcg/mL) due to hyperproduction of beta-lactamase or mutations in penicillin-binding proteins may be misclassified as MSSA on some screening agar; this mechanism (BORSA) is rare.
\end{itemize}
\section*{Risks}
Nasal swabbing is essentially risk-free; it may cause minor discomfort, tickling sensation, or brief nosebleed (epistaxis). Nasopharyngeal swabbing may cause transient discomfort, gagging, sneezing, coughing, or tearing. Rarely (<1\%), nasal swabbing may cause vasovagal syncope.

\section*{What Happens After the Test}
MRSA screening by culture: Results available in 24--48 hours (negative at 24 hours, positive typically at 18--24 hours). MRSA screening by PCR: Results available in 1--2 hours. Respiratory viral PCR: Results available in 1--6 hours depending on platform. Patients identified as MRSA carriers should be placed in contact precautions and considered for decolonization with intranasal mupirocin twice daily for 5 days plus chlorhexidine bathing daily for 5 days \cite{huang2013targeted}. \textit{S. aureus} carriers (MSSA or MRSA) undergoing high-risk surgery should receive intranasal mupirocin decolonization preoperatively plus preoperative vancomycin (for MRSA) or cefazolin (for MSSA) prophylaxis if protocol-driven.

\section*{Understanding Results}

\subsection*{Below Normal}
\begin{itemize}
  \item \textbf{Negative nasal swab (no MRSA or no \textit{S. aureus} detected):} Indicates the patient is not colonized with MRSA or MSSA. Sensitivity of a single anterior nares swab for detecting \textit{S. aureus} carriage is 80--90\%; sampling additional sites (throat, axillae, perineum) increases sensitivity to 95\%. A negative MRSA PCR has a negative predictive value of 98--99\% for ruling out MRSA pneumonia, influencing de-escalation of empiric antibiotics
  \item A negative nasopharyngeal swab PCR does not exclude the target infection if the specimen was poorly collected (inadequate insertion depth, insufficient rotation) or collected after the peak viral shedding period
\end{itemize}

\subsection*{Above Normal}
\begin{itemize}
  \item \textbf{MRSA detected on nasal screening:} Indicates colonization with methicillin-resistant \textit{S. aureus}. Approximately 1--3\% of the general population and 5--15\% of hospitalized patients are MRSA carriers. Colonization is a major risk factor for subsequent MRSA infection (surgical site infection, bloodstream infection, pneumonia), with a 10--30\% risk of developing infection over 18 months. MRSA identification guides vancomycin empiric therapy if infection develops and prompts contact isolation
  \item \textbf{MSSA detected on preoperative screen:} Indicates actionable \textit{S. aureus} carriage. Decolonization with intranasal mupirocin reduces post-operative \textit{S. aureus} infections by approximately 50\%. Vancomycin surgical prophylaxis should not be substituted for cefazolin in MSSA carriers (cefazolin is superior for preventing MSSA SSIs)
  \item \textbf{Respiratory pathogen detected on NP swab:} Confirms active infection with the identified virus or \textit{B. pertussis}. Correlation with symptoms and timing is essential; some viruses (rhinovirus, adenovirus) may be detected for weeks after resolution of acute illness
\end{itemize}

\section*{Normal Results}
\textbf{No MRSA isolated} or \textbf{MRSA not detected} (by PCR). For respiratory viruses: \textbf{Negative} or \textbf{Not detected}. For \textit{B. pertussis}: \textbf{Not detected} by PCR; \textbf{No growth} on culture.

\section*{Follow-up Tests}
\begin{itemize}
  \item \textbf{Blood cultures and site-specific cultures:} When a patient known to be MRSA-colonized develops clinical signs of infection (fever, leukocytosis, purulence); empiric vancomycin should be initiated pending culture results
  \item \textbf{Repeat MRSA screen:} May be warranted after decolonization protocol to document clearance in select populations (recurrent infections); however, recolonization occurs in 30--50\% within weeks to months
  \item \textbf{Confirmatory respiratory pathogen testing:} If rapid antigen testing from NP swab is negative but clinical suspicion for influenza, RSV, or SARS-CoV-2 is high, RT-PCR should be performed due to higher sensitivity
  \item \textit{\textbf{S. aureus}} \textbf{antimicrobial susceptibility testing:} For any detected \textit{S. aureus} isolate to determine oxacillin/methicillin susceptibility and vancomycin MIC, which is critical for empiric and definitive antibiotic selection
  \item \textbf{Screen for other MDROs:} In patients identified as MRSA carriers in endemic or outbreak settings, screening for VRE and carbapenem-resistant organisms may be performed per institutional protocol
\end{itemize}

\section*{References}
\bibliographystyle{unsrtnat}
\bibliography{references}

\clearpage
\section*{Regulatory Status and Authorization Date}
This chapter describes a test category, not a specific commercial product. FDA, CDSCO, EU IVDR, and MHRA approval or registration dates therefore cannot be assigned without the assay manufacturer, product name, instrument, and intended use. For laboratory-developed tests, an individual agency approval date may not exist. See the Preface for the interpretation of regulatory status.

\clearpage
